\renewcommand{\thetheorem}{17.b.\arabic{theorem}}
\setcounter{theorem}{0}

% ---------------------------------------------------------------------------
% Provenance of the prose in this file.
%   % Extractor: <model>   the block below was transcribed from Prof. Biran's
%                          handwritten notes by that model
%   % Creator:   <model>   the block below was written by that model and has no
%                          counterpart in the notes
% A tag holds until the next tag.
% ---------------------------------------------------------------------------

% Extractor: Gemini 3.1 Pro
\section{Change of Basis}
\label{sec:change_of_basis}

\begin{question}
How does the representation matrix $[T]$ of a linear map $T: V \to W$ depend on the choices of bases $\mathcal{B}$ and $\mathcal{C}$?
\end{question}

\begin{corollary}[Change of Basis Formula]
% cor:17.b.1
\label{cor:change_of_basis_formula}
Let $V$ and $W$ be finite-di\-men\-sional vector spaces over $K$. Let $n := \dim V$ and $m := \dim W$. Let $\mathcal{B}$ and $\mathcal{B}'$ be two bases for $V$, and let $\mathcal{C}$ and $\mathcal{C}'$ be two bases for $W$. Then, $[\id_V]_{\mathcal{B}}^{\mathcal{B}'} \in \GL_n(K)$ and $[\id_W]_{\mathcal{C}}^{\mathcal{C}'} \in \GL_m(K)$, with inverses given by:
\begin{equation*}
    [\id_V]_{\mathcal{B}}^{\mathcal{B}'} = ([\id_V]_{\mathcal{B}'}^{\mathcal{B}})^{-1} \quad \text{and} \quad [\id_W]_{\mathcal{C}}^{\mathcal{C}'} = ([\id_W]_{\mathcal{C}'}^{\mathcal{C}})^{-1}
\end{equation*}
Let $T: V \to W$ be a linear map. Then, its representation matrix changes according to the formula:
\begin{equation*}
    [T]_{\mathcal{C}'}^{\mathcal{B}'} = [\id_W]_{\mathcal{C}'}^{\mathcal{C}} \cdot [T]_{\mathcal{C}}^{\mathcal{B}} \cdot [\id_V]_{\mathcal{B}}^{\mathcal{B}'}
\end{equation*}
\end{corollary}

\begin{proof}
We can write $T$ as the composition $T = \id_W \circ T \circ \id_V = \id_W \circ (T \circ \id_V)$. Passing to the matrix representations, this yields:
\begin{equation*}
    [T]_{\mathcal{C}'}^{\mathcal{B}'} = [\id_W]_{\mathcal{C}'}^{\mathcal{C}} \cdot [T \circ \id_V]_{\mathcal{C}}^{\mathcal{B}'} = [\id_W]_{\mathcal{C}'}^{\mathcal{C}} \cdot [T]_{\mathcal{C}}^{\mathcal{B}} \cdot [\id_V]_{\mathcal{B}}^{\mathcal{B}'}
\end{equation*}
Also, we have the identity $\id_V = \id_V \circ \id_V$. This implies that:
\begin{equation*}
    I_n = [\id_V]_{\mathcal{B}}^{\mathcal{B}} = [\id_V]_{\mathcal{B}}^{\mathcal{B}'} \cdot [\id_V]_{\mathcal{B}'}^{\mathcal{B}}
\end{equation*}
and similarly:
\begin{equation*}
    I_n = [\id_V]_{\mathcal{B}'}^{\mathcal{B}'} = [\id_V]_{\mathcal{B}'}^{\mathcal{B}} \cdot [\id_V]_{\mathcal{B}}^{\mathcal{B}'}
\end{equation*}
This demonstrates that $[\id_V]_{\mathcal{B}}^{\mathcal{B}'}$ is invertible and that its inverse is indeed $[\id_V]_{\mathcal{B}'}^{\mathcal{B}}$, which is the asserted identity $[\id_V]_{\mathcal{B}}^{\mathcal{B}'} = ([\id_V]_{\mathcal{B}'}^{\mathcal{B}})^{-1}$. The proof for $[\id_W]_{\mathcal{C}}^{\mathcal{C}'}$ follows the exact same logic.
\end{proof}

\begin{definition}[Change of Basis Matrix]
% def:17.b.2
\label{def:change_of_basis_matrix}
Let $V$ be a finite-di\-men\-sional vector space over $K$, and let $n := \dim V$. Let $\mathcal{B}$ and $\mathcal{B}'$ be two bases of $V$. The matrix $[\id_V]_{\mathcal{B}'}^{\mathcal{B}} \in \GL_n(K)$ is called the \newterm{change of basis matrix} between $\mathcal{B}$ and $\mathcal{B}'$, or the \newterm{transition matrix} from basis $\mathcal{B}$ to basis $\mathcal{B}'$.
\end{definition}

\begin{remark}
If $\mathcal{B} = (w_1, \dots, w_n)$ and $\mathcal{B}' = (w_1', \dots, w_n')$, then the transition matrix is explicitly constructed by placing the coordinate vectors of the old basis elements with respect to the new basis into the columns:
\begin{equation*}
    [\id_V]_{\mathcal{B}'}^{\mathcal{B}} = \begin{pmatrix} | & & | \\ [w_1]_{\mathcal{B}'} & \dots & [w_n]_{\mathcal{B}'} \\ | & & | \end{pmatrix}
\end{equation*}
\end{remark}

\begin{proposition}[Representation of Isomorphisms]
% prop:17.b.3
\label{prop:representation_of_isomorphisms}
Let $T: V \to W$ be an isomorphism between two finite-di\-men\-sional vector spaces. Let $\mathcal{B}$ and $\mathcal{C}$ be bases for $V$ and $W$, respectively. Then, the matrix $[T]_{\mathcal{C}}^{\mathcal{B}}$ is an invertible matrix, and furthermore:
\begin{equation*}
    [T^{-1}]_{\mathcal{B}}^{\mathcal{C}} = ([T]_{\mathcal{C}}^{\mathcal{B}})^{-1}
\end{equation*}
\end{proposition}

\begin{proof}
Let $n := \dim V = \dim W$, the two dimensions agreeing because an isomorphism carries a basis to a basis. By the properties of composed linear maps and representation matrices (the composition rule on \cpageref{prop:representation_matrix_of_composition}, together with the Proposition on \cpageref{prop:identity_represents_id} for the last equality), we have:
\begin{equation*}
    [T^{-1}]_{\mathcal{B}}^{\mathcal{C}} \cdot [T]_{\mathcal{C}}^{\mathcal{B}} = [\id_V]_{\mathcal{B}}^{\mathcal{B}} = I_n
\end{equation*}
and likewise:
\begin{equation*}
    [T]_{\mathcal{C}}^{\mathcal{B}} \cdot [T^{-1}]_{\mathcal{B}}^{\mathcal{C}} = [\id_W]_{\mathcal{C}}^{\mathcal{C}} = I_n
\end{equation*}
This proves the claim.
\end{proof}

\subsection{Equivalence and Similarity of Matrices}

\begin{question}
Suppose $T: V \to V$ is an endomorphism. Let $\mathcal{B}$ and $\mathcal{B}'$ be two bases for $V$. What is the relation between $[T]_{\mathcal{B}'}^{\mathcal{B}'}$ and $[T]_{\mathcal{B}}^{\mathcal{B}}$?
\end{question}

\begin{corollary}[Similarity of Representation Matrices]
% cor:17.b.4
\label{cor:similarity_of_representation_matrices}
Let $V$ be a finite-di\-men\-sional vector space, and let $T: V \to V$ be a linear map. Let $\mathcal{B}$ and $\mathcal{B}'$ be two bases for $V$. Then:
\begin{equation*}
    [T]_{\mathcal{B}'}^{\mathcal{B}'} = [\id_V]_{\mathcal{B}'}^{\mathcal{B}} \cdot [T]_{\mathcal{B}}^{\mathcal{B}} \cdot [\id_V]_{\mathcal{B}}^{\mathcal{B}'} = ([\id_V]_{\mathcal{B}}^{\mathcal{B}'})^{-1} \cdot [T]_{\mathcal{B}}^{\mathcal{B}} \cdot [\id_V]_{\mathcal{B}}^{\mathcal{B}'}
\end{equation*}
\end{corollary}

\begin{definition}[Equivalence and Similarity]
% def:17.b.5
\label{def:equivalence_and_similarity}
We define two types of relations between matrices:
\begin{enumerate}
    \item Let $A, B \in \M_{m \times n}(K)$. We say that $A$ and $B$ are \newterm{equivalent} if there exist matrices $P \in \GL_m(K)$ and $Q \in \GL_n(K)$ such that $B = P A Q$.
    \item Let $A, B \in \M_{n \times n}(K)$. We say that $A$ and $B$ are \newterm{similar} if there exists a matrix $P \in \GL_n(K)$ such that $B = P^{-1} A P$.
\end{enumerate}
\end{definition}

\begin{exercise}[Properties and Characterization of Matrix Equivalence]
\label{exc:equivalence_similarity_properties}
\begin{enumerate}[label=\textbf{(\alph*)}]
    \item Show that \qt{equivalence} between $m \times n$ matrices defines an equivalence relation on $\M_{m \times n}(K)$.
    \item Show that two matrices $A, B \in \M_{m \times n}(K)$ are equivalent if and only if there exist two vector spaces $V$ and $W$ with $\dim V = n$ and $\dim W = m$, two bases $\mathcal{B}$ and $\mathcal{B}'$ of $V$, and two bases $\mathcal{C}$ and $\mathcal{C}'$ of $W$ such that:
    \begin{equation*}
        A = [T]_{\mathcal{C}}^{\mathcal{B}} \quad \text{and} \quad B = [T]_{\mathcal{C}'}^{\mathcal{B}'}
    \end{equation*}
    In other words, $A$ is equivalent to $B$ if and only if $A$ and $B$ are matrix representations of the exact same linear map $T$, but with respect to different choices of bases for $V$ and $W$.
    \item Show that \qt{similarity} between $n \times n$ matrices defines an equivalence relation on $\M_{n \times n}(K)$.
    \item Show that two square matrices $M, N \in \M_{n \times n}(K)$ are similar if and only if there exists an $n$-di\-men\-sional space $V$, an endomorphism $T: V \to V$, and two bases $\mathcal{B}$ and $\mathcal{B}'$ of $V$ such that $M = [T]_{\mathcal{B}}^{\mathcal{B}}$ and $N = [T]_{\mathcal{B}'}^{\mathcal{B}'}$.
\end{enumerate}
\exinfo{This exercise is taken from Prof.~Biran's handwritten lecture notes.}
\end{exercise}

\begin{exercise}[Similarity of Matrix Products]
\label{exc:similarity_of_matrix_products}
Let $A, B \in \M_{n \times n}(K)$.
\begin{enumerate}
    \item Show that if at least one of $A$ or $B$ is invertible, then $AB$ and $BA$ are similar.
    \item Does the conclusion hold in general if neither $A$ nor $B$ is invertible? (Provide a proof or a counterexample).
\end{enumerate}
\exinfo{This exercise is Problem 4 of Problem Sheet 13, Linear Algebra I, Autumn Semester 2025.}
\end{exercise}

\begin{exercise}[Rank Invariance under Similarity and Fixed Vectors]
\label{exc:rank_invariance_similarity_mc}
Let $A \in \M_{n \times n}(K)$ and $P \in \GL_n(K)$. Set
\[
r := \rank(A - I_n) \qquad \text{and} \qquad s := \rank(P A P^{-1} - I_n).
\]
Which of the following statements hold?
\begin{enumerate}
    \item $r = s$.
    \item $r \ne s$.
    \item $r > s$.
    \item $r < s$.
    \item If $r < n$, then there exists a non-zero vector $v \in K^n \setminus \{0\}$ such that $Av = v$.
\end{enumerate}
\exinfo{This exercise is Multiple Choice Question 2 of Problem Sheet 13, Linear Algebra I, Autumn Semester 2025.}
\end{exercise}

\begin{exercise}[Change of Basis in Polynomial Spaces]
\label{exc:change_of_basis_polynomials_mc}
Consider the two ordered bases of $\mathbb{R}[x]_2$:
\[
\mathcal{B} = (1, x, x^2), \qquad \mathcal{C} = \left( x^2, (x+1)^2, (x+2)^2 \right).
\]
Determine the change of basis (transition) matrix $Q := [\id_{\mathbb{R}[x]_2}]_{\mathcal{B}}^{\mathcal{C}}$.
\exinfo{This exercise is Multiple Choice Question 1 of Problem Sheet 13, Linear Algebra I, Autumn Semester 2025.}
\end{exercise}

\begin{exercise}[Properties of Linear Maps and Base Changes]
\label{exc:linear_maps_and_base_changes_sc}
Which of the following statements is \emph{false} in general?
\begin{enumerate}
    \item The change of basis matrix is the representation matrix of the identity map with respect to the corresponding bases.
    \item Every finite-dimensional vector space over $K$ is isomorphic to $K^n$ for some $n \ge 0$.
    \item The rank of any linear map $f \colon K^n \to K^m$ is at least $\min\{n, m\}$.
    \item The representation matrix of an isomorphism with respect to any choices of bases is invertible.
\end{enumerate}
\exinfo{This exercise is Single Choice Question 1 of Problem Sheet 13, Linear Algebra I, Autumn Semester 2025.}
\end{exercise}

\begin{proposition}[Normal Form for Linear Maps]
% prop:17.b.6
\label{prop:linear_map_normal_form}
Let $V$ and $W$ be finite-di\-men\-sional vector spaces over $K$, with dimensions $n = \dim V$ and $m = \dim W$. Let $T: V \to W$ be a linear map with $\rank(T) = r$. (Recall that $\rank(T) := \dim \Img(T)$). Then, there exists a basis $\mathcal{B} = (v_1, \dots, v_n)$ of $V$ and a basis $\mathcal{C} = (w_1, \dots, w_m)$ of $W$ such that:
\begin{equation*}
    [T]_{\mathcal{C}}^{\mathcal{B}} = \begin{pmatrix} I_r & O_{r, n-r} \\ O_{m-r, r} & O_{m-r, n-r} \end{pmatrix}
\end{equation*}
where $O_{p,q}$ stands for the zero matrix in $\M_{p \times q}(K)$.
\end{proposition}

\begin{proof}
Put $\ell := \dim \ker(T)$. By the rank theorem, we know that $r = n - \ell$. Let $v_1, \dots, v_\ell$ be a basis of $\ker(T)$, and extend it to a full basis $(v_1, \dots, v_\ell, v_{\ell+1}, \dots, v_n)$ of $V$. It will be useful to work below with the rearranged basis:
\begin{equation*}
    \mathcal{B} := (v_{\ell+1}, \dots, v_n, v_1, \dots, v_\ell)
\end{equation*}
In the proof of the rank theorem, we have seen that $T(v_{\ell+1}), \dots, T(v_n)$ form a basis for $\Img(T)$. Define $w_i := T(v_{\ell+i})$ for $i = 1, \dots, r$ (i.e., $w_1 = T(v_{\ell+1}), \dots, w_r = T(v_{\ell+r})$), and extend this list to a basis of $W$:
\begin{equation*}
    \mathcal{C} := (w_1, \dots, w_r, w_{r+1}, \dots, w_m)
\end{equation*}
% Creator: Opus 5
It now directly follows from these definitions that the representation matrix $[T]_{\mathcal{C}}^{\mathcal{B}}$ takes the exact block form described in the proposition. Indeed, the $j$-th column of $[T]_{\mathcal{C}}^{\mathcal{B}}$ is the coordinate vector of the image of the $j$-th member of $\mathcal{B}$. For $1 \le j \le r$ that member is $v_{\ell+j}$, whose image is $w_j$ by construction, so the column is $[w_j]_{\mathcal{C}} = e_j \in K^m$. For $r < j \le n$ the member is one of $v_1, \dots, v_\ell$, which lie in $\ker(T)$, so the column is $[0_W]_{\mathcal{C}} = 0$. A matrix whose first $r$ columns are $e_1, \dots, e_r$ and whose remaining columns vanish is exactly the asserted block matrix.
% Extractor: Gemini 3.1 Pro
\end{proof}

\begin{corollary}[Matrix Equivalence and Normal Form]
% cor:17.b.7
\label{cor:matrix_equivalence_normal_form}
Let $A \in \M_{m \times n}(K)$. Then, there exist invertible matrices $P \in \GL_m(K)$ and $Q \in \GL_n(K)$ such that:
\begin{equation*}
    PAQ = \begin{pmatrix} I_r & O \\ O & O \end{pmatrix}
\end{equation*}
where $r = \rankcol(A)$. In particular, $\rankcol(A) \le \min\{m, n\}$.
\end{corollary}

\begin{proof}
Consider the linear map $T_A: K^n \to K^m$ defined by left-multiplication, $T_A(v) = A \cdot v$. Then, $r = \rankcol(A) = \dim \Img(T_A)$, because we have previously seen that 
\[\Img(T_A) = \Sp(T_A(e_1), \dots, T_A(e_n)) = \Sp(\text{columns of } A).\]

Now, we apply the previous proposition to $T_A$, yielding that:
\begin{equation*}
    [T_A]_{\mathcal{C}}^{\mathcal{B}} = \begin{pmatrix} I_r & O \\ O & O \end{pmatrix}
\end{equation*}
for some basis $\mathcal{B}$ of $K^n$ and some basis $\mathcal{C}$ of $K^m$. Using the change of basis formula, we can rewrite this as:
\begin{equation*}
    \underbrace{[\id_{K^m}]_{\mathcal{C}}^{\mathcal{E}_m}}_{=: P} \cdot \underbrace{[T_A]_{\mathcal{E}_m}^{\mathcal{E}_n}}_{= A} \cdot \underbrace{[\id_{K^n}]_{\mathcal{E}_n}^{\mathcal{B}}}_{=: Q} = \begin{pmatrix} I_r & O \\ O & O \end{pmatrix}
\end{equation*}
where $\mathcal{E}_n$ and $\mathcal{E}_m$ are the standard bases of $K^n$ and $K^m$, respectively, and where the bases meet diagonally as they must: the superscript $\mathcal{E}_m$ of the left factor matches the subscript $\mathcal{E}_m$ of the middle one, and the subscript $\mathcal{E}_n$ of the right factor matches the superscript $\mathcal{E}_n$ of the middle one. Setting $P := [\id_{K^m}]_{\mathcal{C}}^{\mathcal{E}_m}$ and $Q := [\id_{K^n}]_{\mathcal{E}_n}^{\mathcal{B}}$ concludes the proof.
\end{proof}

\begin{remark}
The fact that $\rankcol(A) \le \min\{m, n\}$ also follows from observing the linear map $T_A: K^n \to K^m$. We have:
\begin{equation*}
    \rankcol(A) = \dim \Img(T_A) = n - \dim \ker(T_A) \le n
\end{equation*}
and simultaneously,
\begin{equation*}
    \rankcol(A) = \dim \Img(T_A) \le \dim K^m = m
\end{equation*}
\end{remark}

\begin{question}
Recall that we have an equivalence relation on $\M_{m \times n}(K)$: $A \sim B$ if there exist $P \in \GL_m(K)$ and $Q \in \GL_n(K)$ such that $B = PAQ$. Can we describe the equivalence classes of this equivalence relation?
\end{question}

\begin{corollary}[Classification of Matrices by Rank]
% cor:17.b.8
\label{cor:matrix_rank_equivalence_classification}
Let $A, B \in \M_{m \times n}(K)$. Then, $A \sim B$ if and only if $\rankcol(A) = \rankcol(B)$. Moreover, for every $0 \le r \le \min\{m, n\}$, there is precisely one equivalence class of matrices $A$ with $\rankcol = r$.
\end{corollary}

\begin{proof}
Let $A, B \in \M_{m \times n}(K)$. First, assume that $\rankcol(A) = \rankcol(B)$, and denote this common column rank by $r$. By \cref{cor:matrix_equivalence_normal_form}, we know that $A \sim \begin{pmatrix} I_r & 0 \\ 0 & 0 \end{pmatrix}$ and $B \sim \begin{pmatrix} I_r & 0 \\ 0 & 0 \end{pmatrix}$. Since $\sim$ is an equivalence relation, it naturally follows that $A \sim B$.

Conversely, assume that $A \sim B$. This implies that there exist matrices $P \in \GL_m(K)$ and $Q \in \GL_n(K)$ such that $B = PAQ$. Transitioning to the associated linear maps, this means $T_B = T_P \circ T_A \circ T_Q$. We analyze the image of $T_B$:
\begin{equation*}
    \Img(T_B) = T_B(K^n) = T_P \circ T_A(T_Q(K^n)) = T_P \circ T_A(K^n) = T_P(\Img(T_A))
\end{equation*}
Note that $T_Q(K^n) = K^n$ because $T_Q$ is an isomorphism. Consequently, the dimension of the image is:
\begin{equation*}
    \rankcol(B) = \dim \Img(T_B) = \dim T_P(\Img(T_A)) = \dim \Img(T_A) = \rankcol(A)
\end{equation*}
where the second to last equality holds because $T_P$ is an isomorphism, preserving the dimension of subspaces.

% Creator: Opus 5
The last statement of the corollary now follows. The two implications just proved say that the equivalence classes are exactly the sets of matrices of a common column rank, so distinct ranks give distinct classes, and every value $0 \le r \le \min\{m, n\}$ actually occurs, since
\[
\begin{pmatrix} I_r & O \\ O & O \end{pmatrix} \in \M_{m \times n}(K)
\]
has columns $e_1, \dots, e_r, 0, \dots, 0$ and hence column rank $r$. There is therefore precisely one class for each such $r$.
\end{proof}

% Creator: Opus 5
\subsection{Solutions to the Exercises}
\label{sec:solutions_change_of_basis}

\begin{exercisesolution}[Proof of \cref{exc:equivalence_similarity_properties}]
Throughout, write $A \sim B$ for equivalence and $A \approx B$ for similarity.

\begin{enumerate}[label=\textbf{(\alph*)}]
    \item \textbf{Equivalence is an equivalence relation.} It is reflexive, since $A = I_m A I_n$ with $I_m \in \GL_m(K)$ and $I_n \in \GL_n(K)$. It is symmetric: if $B = PAQ$ with $P \in \GL_m(K)$ and $Q \in \GL_n(K)$, then multiplying by $P^{-1}$ on the left and $Q^{-1}$ on the right gives $A = P^{-1} B Q^{-1}$, and $P^{-1}, Q^{-1}$ are again invertible by \cref{prop:gl_group_structure}. It is transitive: if $B = PAQ$ and $C = P'BQ'$, then
    \[
    C = P'(PAQ)Q' = (P'P) A (QQ'),
    \]
    and $P'P \in \GL_m(K)$, $QQ' \in \GL_n(K)$, again by \cref{prop:gl_group_structure}.

    \item \textbf{Equivalent matrices represent one map in two pairs of bases.} Suppose first that $A = [T]_{\mathcal{C}}^{\mathcal{B}}$ and $B = [T]_{\mathcal{C}'}^{\mathcal{B}'}$ for a single linear map $T: V \to W$. By \cref{cor:change_of_basis_formula},
    \[
    B = \underbrace{[\id_W]_{\mathcal{C}'}^{\mathcal{C}}}_{\in \, \GL_m(K)} \cdot A \cdot \underbrace{[\id_V]_{\mathcal{B}}^{\mathcal{B}'}}_{\in \, \GL_n(K)},
    \]
    so $A \sim B$.

    Conversely, suppose $B = PAQ$ with $P \in \GL_m(K)$ and $Q \in \GL_n(K)$. Take $V := K^n$, $W := K^m$ and $T := T_A$, and let $\mathcal{E}_n$, $\mathcal{E}_m$ be the standard bases, so that $A = [T]_{\mathcal{E}_m}^{\mathcal{E}_n}$. Let $\mathcal{B}'$ be the list of columns of $Q$ and let $\mathcal{C}'$ be the list of columns of $P^{-1}$; both are bases, since a square matrix is invertible precisely when its columns form a basis. The matrix $[\id]_{\mathcal{E}_n}^{\mathcal{B}'}$ has as its $j$-th column the coordinate vector of the $j$-th member of $\mathcal{B}'$ with respect to the standard basis, that is, the $j$-th column of $Q$; hence $[\id_{K^n}]_{\mathcal{E}_n}^{\mathcal{B}'} = Q$, and likewise $[\id_{K^m}]_{\mathcal{E}_m}^{\mathcal{C}'} = P^{-1}$, so that $[\id_{K^m}]_{\mathcal{C}'}^{\mathcal{E}_m} = P$ by \cref{cor:change_of_basis_formula}. That same corollary now gives
    \[
    [T]_{\mathcal{C}'}^{\mathcal{B}'} = [\id_{K^m}]_{\mathcal{C}'}^{\mathcal{E}_m} \cdot [T]_{\mathcal{E}_m}^{\mathcal{E}_n} \cdot [\id_{K^n}]_{\mathcal{E}_n}^{\mathcal{B}'} = P A Q = B,
    \]
    so $A$ and $B$ do represent the same map in two pairs of bases.

    \item \textbf{Similarity is an equivalence relation.} Reflexivity: $A = I_n^{-1} A I_n$. Symmetry: if $B = P^{-1}AP$, then $A = P B P^{-1} = (P^{-1})^{-1} B (P^{-1})$. Transitivity: if $B = P^{-1}AP$ and $C = R^{-1}BR$, then
    \[
    C = R^{-1} P^{-1} A P R = (PR)^{-1} A (PR),
    \]
    using $(PR)^{-1} = R^{-1}P^{-1}$ from \cref{prop:gl_group_structure}.

    \item \textbf{Similar matrices represent one endomorphism in two bases.} If $M = [T]_{\mathcal{B}}^{\mathcal{B}}$ and $N = [T]_{\mathcal{B}'}^{\mathcal{B}'}$ for an endomorphism $T: V \to V$, then \cref{cor:similarity_of_representation_matrices} gives $N = P^{-1} M P$ with $P := [\id_V]_{\mathcal{B}}^{\mathcal{B}'} \in \GL_n(K)$, so $M \approx N$.

    Conversely, let $N = P^{-1} M P$ with $P \in \GL_n(K)$. Take $V := K^n$, $T := T_M$ and $\mathcal{B} := \mathcal{E}_n$, so that $M = [T]_{\mathcal{B}}^{\mathcal{B}}$. Let $\mathcal{B}'$ be the list of columns of $P$, a basis of $K^n$, so that $[\id_{K^n}]_{\mathcal{E}_n}^{\mathcal{B}'} = P$ as in \textbf{(b)}. Then \cref{cor:similarity_of_representation_matrices} gives
    \[
    [T]_{\mathcal{B}'}^{\mathcal{B}'} = \big([\id_{K^n}]_{\mathcal{E}_n}^{\mathcal{B}'}\big)^{-1} \cdot [T]_{\mathcal{E}_n}^{\mathcal{E}_n} \cdot [\id_{K^n}]_{\mathcal{E}_n}^{\mathcal{B}'} = P^{-1} M P = N. \qedhere
    \]
\end{enumerate}
\end{exercisesolution}

\begin{ainote}
Every numbered statement of the notes for this section is present in the
transcript, in the notes' order, and 17.b.1--17.b.8 already matched the
handwritten numbering; the proofs follow Prof. Biran's.

The four parts of \cref{exc:equivalence_similarity_properties} are Prof. Biran's,
marked \qt{Exc.} in the notes and left unanswered by the transcript; they are
solved above. Two of them carry the content of the section: equivalence of
matrices means representing the same linear map in two pairs of bases, and
similarity means representing the same endomorphism in two bases. Both
directions turn on the observation that for an invertible $P$, the columns of
$P$ form a basis whose transition matrix from the standard basis is $P$ itself.

Two slips in the bookkeeping of bases were corrected, and the second is the kind
this book is most exposed to. In the proof of
\cref{cor:matrix_equivalence_normal_form} the left factor was written
$[\id_{K^m}]_{\mathcal{E}_m}^{\mathcal{C}}$, whose bases do not meet those of the
middle factor $[T_A]_{\mathcal{E}_m}^{\mathcal{E}_n}$ the way the composition
rule requires: the product as printed does not compose at all, and the matrix
named $P$ was the inverse of the one intended. It has to be
$[\id_{K^m}]_{\mathcal{C}}^{\mathcal{E}_m}$, which is also what the solution to
part \textbf{(b)} of \cref{exc:equivalence_similarity_properties} arrives at
independently. The other slip was one inverse too many at the close of
\cref{cor:change_of_basis_formula}.
\end{ainote}